\input{preamble}

% OK, start here.
%
\begin{document}

\title{Introduction aux champs algébriques}

\maketitle

\phantomsection
\label{section-phantom}

\tableofcontents




\section{Pourquoi lire ce chapitre ?}
\label{section-introduction}

\noindent
Nous donnons une introduction informelle aux champs algébriques. Le but est
d'introduire rapidement un langage simple qui permette de penser les
propriétés locales et globales de votre problème de modules favori.
Une fois cela fait, vous devriez pouvoir formuler des questions bien posées
sur les problèmes de modules et commencer à les résoudre, en admettant
l'existence d'une théorie générale. Si vous aboutissez à un résultat
intéressant, vous pourrez revenir à la théorie générale développée dans les
autres parties du projet Champs et combler les lacunes selon les besoins.

\medskip\noindent
Le point de vue adopté ici est proche de celui de
\cite{KatzMazur} et \cite{mumford_picard}.






\section{Préliminaires}
\label{section-preliminary}

\noindent
Soit $S$ un schéma. Une {\it courbe elliptique} sur $S$ est un triplet
$(E, f, 0)$, où $E$ est un schéma et où $f : E \to S$ et $0 : S \to E$
sont des morphismes de schémas tels que
\begin{enumerate}
\item $f : E \to S$ soit propre et lisse de dimension relative $1$ ;
\item pour tout $s \in S$, la fibre $E_s$ soit une courbe connexe
de genre $1$, c'est-à-dire que $H^0(E_s, \mathcal{O})$ et
$H^1(E_s, \mathcal{O})$ soient tous deux des espaces vectoriels de dimension
$1$ sur $\kappa(s)$ ;
\item $0$ soit une section de $f$.
\end{enumerate}
Étant données des courbes elliptiques $(E, f, 0)/S$ et $(E', f', 0')/S'$, un
{\it morphisme de courbes elliptiques au-dessus de $a : S \to S'$}
est un morphisme $\alpha : E \to E'$ tel que le diagramme
$$
\xymatrix{
E \ar[rr]_\alpha \ar[d]^f & & E' \ar[d]_{f'}  \\
S \ar@/^5ex/[u]^0 \ar[rr]^a & & S' \ar@/_5ex/[u]_{0'}
}
$$
soit commutatif et que le carré intérieur soit cartésien ; autrement dit, le
morphisme $\alpha$ induit un isomorphisme $E \to S \times_{S'} E'$.
Nous allons définir le champ des courbes elliptiques $\mathcal{M}_{1, 1}$.
Dans le reste du projet Champs, nous développons la méthode introduite dans
l'article de Deligne et Mumford \cite{DM}, qui consiste à présenter
$\mathcal{M}_{1, 1}$ comme une catégorie munie d'un foncteur
$$
p : \mathcal{M}_{1, 1} \longrightarrow \Sch, \quad
(E, f, 0)/S \longmapsto S
$$
Cela conduit à travailler avec les catégories fibrées sur la catégorie des
schémas, les topologies, les champs fibrés en groupoïdes, les recouvrements,
etc. Dans ce chapitre, nous laissons tout cet appareil de côté et envisageons
les choses un peu autrement, probablement d'une manière plus proche de celle
dont les fondateurs de la théorie ont commencé à les concevoir.


\section{Le champ de modules des courbes elliptiques}
\label{section-moduli-elliptic-curves}

\noindent
Voici la démarche que nous allons suivre :
\begin{enumerate}
\item Partons de votre catégorie de schémas favorite $\Sch$.
\item Ajoutons un nouveau symbole $\mathcal{M}_{1, 1}$.
\item Un morphisme $S \to \mathcal{M}_{1, 1}$ {\bf est} une courbe elliptique
$(E, f, 0)$ sur $S$.
\item Un diagramme
$$
\xymatrix{
S \ar[rr]_a \ar[rd]_{(E, f, 0)} & & S' \ar[ld]^{(E', F', 0')} \\
& \mathcal{M}_{1, 1}
}
$$
{\bf est} commutatif si et seulement s'il existe un morphisme
$\alpha : E \to E'$ de courbes elliptiques au-dessus de $a : S \to S'$.
Nous disons que $\alpha$ {\it atteste} la commutativité du diagramme.
\item Remarquons que les diagrammes commutatifs se composent comme suit :
$$
\xymatrix{
S \ar[rrr]_a \ar[rrrd]_{(E, f, 0)} & & &
S' \ar[d]_{(E', F', 0')} \ar[rrr]_{a'} & & &
S'' \ar[llld]^{(E'', F'', 0'')}
\\
& & & \mathcal{M}_{1, 1}
}
$$
en effet, $\alpha' \circ \alpha$ atteste la commutativité du triangle extérieur
si $\alpha$ et $\alpha'$ attestent respectivement la commutativité des triangles
de gauche et de droite.
\item La composée
$$
S \xrightarrow{a} S' \xrightarrow{(E', f', 0')} \mathcal{M}_{1, 1}
$$
est donnée par $(E' \times_{S'} S, f' \times_{S'} S, 0' \times_{S'} S)$.
\end{enumerate}
Au terme de cette procédure, nous avons agrandi la catégorie $\Sch$ des
schémas en lui adjoignant exactement un objet...

\medskip\noindent
À ceci près que nous n'avons pas défini ce qu'est un morphisme de
$\mathcal{M}_{1, 1}$ vers un schéma $T$. La réponse consiste à retenir la
notion la plus faible possible pour laquelle les compositions aient un sens.
Ainsi, un morphisme $F : \mathcal{M}_{1, 1} \to T$ est une règle qui associe à
toute courbe elliptique $(E, f, 0)/S$ un morphisme
$F(E, f, 0) : S \to T$, de telle sorte que, pour tout diagramme commutatif
$$
\xymatrix{
S \ar[rr]_a \ar[rd]_{(E, f, 0)} & & S' \ar[ld]^{(E', F', 0')} \\
& \mathcal{M}_{1, 1}
}
$$
le diagramme
$$
\xymatrix{
S \ar[rr]_a \ar[rd]_{F(E, f, 0)} & & S' \ar[ld]^{F(E', F', 0')} \\
& T
}
$$
soit lui aussi commutatif. Le $j$-invariant, dont vous avez peut-être entendu
parler, en fournit un exemple :
$$
j : \mathcal{M}_{1, 1} \longrightarrow \mathbf{A}^1_{\mathbf{Z}}
$$
Bien, nous avons donc terminé...

\medskip\noindent
Eh bien non ! Il nous faut encore définir une notion de morphisme
$\mathcal{M}_{1, 1} \to \mathcal{M}_{1, 1}$. Nous procédons exactement comme
précédemment : un morphisme
$F : \mathcal{M}_{1, 1} \to \mathcal{M}_{1, 1}$
est une règle qui associe à toute courbe elliptique $(E, f, 0)/S$ une autre
courbe elliptique $F(E, f, 0)$, tout en préservant la commutativité des
diagrammes précédents. Toutefois, comme je ne connais aucun exemple non
trivial d'un tel foncteur, je définirai pour l'instant l'ensemble des
morphismes de $\mathcal{M}_{1, 1}$ vers lui-même comme réduit à l'identité.

\medskip\noindent
J'espère que vous voyez comment ajouter d'autres objets à cette catégorie
agrandie. Il semble intuitivement clair que, pour tout problème de modules
« bien sage », nous pouvons effectuer la construction précédente et ajouter
un objet à notre catégorie. En fait, une grande partie de la géométrie
algébrique contemporaine se déroule dans un tel univers, où $\Sch$ est
agrandie par une infinité dénombrable de champs de modules construits
explicitement.

\medskip\noindent
Vous pourriez objecter que la structure obtenue n'est pas une catégorie,
car il subsiste un certain « flou » quant aux diagrammes qui commutent et aux
combinaisons de diagrammes qui continuent de commuter : il faut en effet
produire une attestation de la commutativité. Or il se trouve que cette idée,
qui consiste à munir la commutativité d'attestations, constitue une approche
valable des $2$-catégories ! Nous nous y tiendrons donc.






\section{Produits fibrés}
\label{section-fibre-products}

\noindent
La question posée ici est de déterminer ce que doit être le produit fibré
$$
\xymatrix{
& ? \ar@{..>}[rd] \ar@{..>}[ld] \\
S \ar[rd]_{(E, f, 0)} & & S' \ar[ld]^{(E', f', 0')} \\
& \mathcal{M}_{1, 1}
}
$$
La réponse est la suivante : un morphisme d'un schéma $T$ vers $?$
doit être un triplet $(a, a', \alpha)$, où
$a : T \to S$ et $a' : T \to S'$ sont des morphismes de schémas et où
$\alpha : E \times_{S, a} T \to E' \times_{S', a'} T$ est un isomorphisme de
courbes elliptiques sur $T$. Cela a un sens en vertu de notre définition
antérieure de la composition et des diagrammes commutatifs.

\begin{lemma}[Fait clé]
\label{lemma-key-fact}
Le foncteur $\Sch^{opp} \to \textit{Ensembles}$,
$T \mapsto \{(a, a', \alpha)\text{ comme ci-dessus}\}$
est représentable par un schéma $S \times_{\mathcal{M}_{1, 1}} S'$.
\end{lemma}

\begin{proof}
Idée de la démonstration. Relier ce foncteur à
$$
\mathit{Isom}_{S \times S'}(E \times S', S \times E')
$$
et utiliser la théorie de Grothendieck des schémas de Hilbert.
\end{proof}

\begin{remark}
\label{remark-diagonal}
Nous avons la formule
$S \times_{\mathcal{M}_{1, 1}} S' =
(S \times S')
\times_{\mathcal{M}_{1, 1} \times \mathcal{M}_{1, 1}}
\mathcal{M}_{1, 1}$.
Le fait clé est donc une propriété de la diagonale
$\Delta_{\mathcal{M}_{1, 1}}$ de $\mathcal{M}_{1, 1}$.
\end{remark}

\noindent
Quoi qu'il en soit, le fait clé nous permet de poser la définition suivante.

\begin{definition}
\label{definition-smooth}
Nous disons qu'un morphisme $S \to \mathcal{M}_{1, 1}$ est {\it lisse} si,
pour tout morphisme $S' \to \mathcal{M}_{1, 1}$, le morphisme de projection
$$
S \times_{\mathcal{M}_{1, 1}} S' \longrightarrow S'
$$
est lisse.
\end{definition}

\noindent
Remarquons que cette définition est compatible avec la notion de morphisme
lisse de schémas, puisque tout changement de base d'un morphisme lisse est
lisse. Il est en outre clair comment étendre cette définition à d'autres
propriétés des morphismes vers $\mathcal{M}_{1, 1}$ (ou vers votre champ de
modules favori). Nous l'utiliserons notamment ci-dessous pour les morphismes
{\it surjectifs}.





\section{La définition}
\label{section-definition}

\noindent
Nous allons la formuler comme une définition plutôt que comme un résultat,
car nous attendons du lecteur qu'il essaie d'autres cas, et pas seulement le
champ $\mathcal{M}_{1, 1}$ ni la seule catégorie $\Sch$ de tous les schémas.

\begin{definition}
\label{definition-algebraic-stack}
Nous disons que $\mathcal{M}_{1, 1}$ est un {\it champ algébrique} si et
seulement si
\begin{enumerate}
\item les objets vérifient la descente pour la topologie \'etale sur $\Sch$ ;
\item le fait clé est vérifié ;
\item il existe un morphisme surjectif et lisse
$S \to \mathcal{M}_{1, 1}$.
\end{enumerate}
\end{definition}

\noindent
La première condition est une « propriété de faisceau ». Nous allons
l'expliciter, car elle recèle un point technique. Supposons donnés un schéma
$S$, un recouvrement \'etale $\{S_i \to S\}$ et des morphismes
$e_i : S_i \to \mathcal{M}_{1, 1}$ tels que les diagrammes
$$
\xymatrix{
S_i \times_S S_j \ar[rd]_{e_i \circ \text{pr}_1} \ar[rr]_{\text{id}} & &
S_i \times_S S_j \ar[ld]^{e_j \circ \text{pr}_2} \\
& \mathcal{M}_{1, 1}
}
$$
soient commutatifs. La condition de faisceau ne garantit {\it pas}, à elle
seule, l'existence d'un morphisme $e : S \to \mathcal{M}_{1, 1}$ dans cette
situation. En effet, il faut choisir des attestations $\alpha_{ij}$ pour les
diagrammes précédents et imposer que
$$
\text{pr}_{02}^*\alpha_{ik} =
\text{pr}_{12}^*\alpha_{jk} \circ \text{pr}_{01}^*\alpha_{ij}
$$
comme attestations sur $S_i \times_S S_j \times_S S_k$. Le sens de cette
condition devrait être clair... Sinon, je crains qu'il ne vous faille lire
une partie des développements sur les catégories fibrées en groupoïdes, etc.
Quoi qu'il en soit, l'égalité affichée est souvent appelée la
{\it condition de cocycle}. Voici une formulation plus précise de la
« propriété de faisceau » : étant donnés $\{S_i \to S\}$,
$e_i : S_i \to \mathcal{M}_{1, 1}$ et des attestations $\alpha_{ij}$
satisfaisant la condition de cocycle, il existe un unique morphisme
$e : S \to \mathcal{M}_{1, 1}$, à isomorphisme unique près, tel que
$e_i \cong e|_{S_i}$ et qui redonne les $\alpha_{ij}$.

\medskip\noindent
Comme vous le voyez, la formulation même d'un énoncé précis demande déjà un
peu de travail. La démonstration de cette « propriété de faisceau » repose sur
une technique fondamentale de la géométrie algébrique : la théorie de la
descente. Je suggère, dans un premier temps, d'admettre simplement cette
« propriété de faisceau » et d'examiner ce qu'elle implique en pratique.
Une certaine agilité d'esprit est en effet nécessaire pour condenser cette
« propriété de faisceau » en un énoncé maniable tenant sur un coin de table.
La variante la plus simple qui présente déjà quelque intérêt est peut-être la
suivante. Supposons donnée une extension galoisienne finie $L/K$ de corps, de
groupe de Galois $G = \text{Gal}(L/K)$. Posons $T = \Spec(L)$ et
$S = \Spec(K)$. Alors $\{T \to S\}$ est un recouvrement \'etale.
Soit $(E, f, 0)$ une courbe elliptique sur $L$. (Cela signifie simplement que
$E \subset \mathbf{P}^2_L$ est donnée par une équation de Weierstrass et que
$0$ est le point à l'infini usuel.) Notons
$E_\sigma = E \times_{T, \Spec(\sigma)} T$ le changement de base.
(Cela revient à appliquer $\sigma$ aux coefficients de l'équation de
Weierstrass, à moins que ce ne soit $\sigma^{-1}$.) Supposons en outre que,
pour tout $\sigma \in G$, nous soit donné un isomorphisme
$$
\alpha_\sigma : E \longrightarrow E_\sigma
$$
sur $T$. Dans cette situation, la condition de cocycle signifie que
$$
(\alpha_\tau)^\sigma \circ \alpha_\sigma = \alpha_{\tau\sigma}
$$
pour $\sigma, \tau \in G$. Si vous avez déjà pratiqué la cohomologie des
groupes, cette formule vous sera familière. La condition de « recollement »
sur $\mathcal{M}_{1, 1}$ affirme donc que, si ce système d'équations admet une
solution, il existe une courbe elliptique $E'$ sur $S$ telle que
$E \cong E' \times_S T$ (elle en dit un peu plus, puisqu'elle indique aussi
comment retrouver les $\alpha_\sigma$).

\medskip\noindent
Défi : pouvez-vous le démontrer en n'utilisant que des courbes elliptiques
définies par des équations de Weierstrass ?




\section{Un recouvrement lisse}
\label{section-smooth}

\noindent
Il nous reste à trouver un recouvrement lisse de $\mathcal{M}_{1, 1}$.
En un certain sens, l'existence d'un recouvrement lisse
{\it implique}\footnote{Il y a là une légère tricherie : pour vérifier la
lissité, il faut démontrer un résultat proche du fait clé, puisque la lissité
est définie au moyen de produits fibrés. L'avantage est qu'il suffit de
démontrer l'existence de ces produits fibrés lorsque, d'un côté, figure le
morphisme dont on cherche précisément à montrer qu'il fournit le recouvrement
lisse.} le fait clé ! Dans le cas des courbes elliptiques, nous en construisons
un au moyen de l'équation de Weierstrass.

\medskip\noindent
Posons
$$
W = \Spec(\mathbf{Z}[a_1, a_2, a_3, a_4, a_6, 1/\Delta])
$$
où $\Delta \in \mathbf{Z}[a_1, a_2, a_3, a_4, a_6]$ est un certain polynôme
(voir ci-dessous). Posons
$$
\mathbf{P}_W^2 \supset E_W :
zy^2 + a_1 xyz + a_3 yz^2 = x^3 + a_2x^2z + a_4xz^2 + a_6z^3.
$$
Notons $f_W : E_W \to W$ la projection. Enfin, notons $0_W : W \to E_W$ la
section de $f_W$ donnée par $(0 : 1 : 0)$. Il existe alors un polynôme homogène
de degré $12$, noté $\Delta$, en les $a_i$, où $\deg(a_i) = i$, tel que
$E_W \to W$ soit lisse. On peut l'obtenir explicitement en calculant les
dérivées partielles de l'équation de Weierstrass ; on peut bien sûr aussi le
chercher dans la littérature. On peut encore demander à pari/gp de le
calculer. Le voici :
\begin{align*}
\Delta & = -a_6a_1^6 + a_4a_3a_1^5 + ((-a_3^2 - 12a_6)a_2 + a_4^2)a_1^4 + \\
& (8a_4a_3a_2 + (a_3^3 + 36a_6a_3))a_1^3 + \\
& ((-8a_3^2 - 48a_6)a_2^2 + 8a_4^2a_2 + (-30a_4a_3^2 + 72a_6a_4))a_1^2 + \\
& (16a_4a_3a_2^2 + (36a_3^3 + 144a_6a_3)a_2 - 96a_4^2a_3)a_1 + \\
& (-16a_3^2 - 64a_6)a_2^3 + 16a_4^2a_2^2 + (72a_4a_3^2 + 288a_6a_4)a_2 + \\
& -27a_3^4 - 216a_6a_3^2  -64a_4^3 - 432a_6^2
\end{align*}
Vous reconnaîtrez peut-être les deux derniers termes du cas
$y^2 = x^3 + Ax + B$, où le discriminant vaut
$-64A^3 - 432B^2 = -16(4A^3 + 27B^2)$.

\begin{lemma}
\label{lemma-Weierstrass-smooth-cover}
Le morphisme $W \xrightarrow{(E_W, f_W, 0_W)} \mathcal{M}_{1, 1}$ est lisse
et surjectif.
\end{lemma}

\begin{proof}
La surjectivité résulte du fait que toute courbe elliptique sur un corps admet
une équation de Weierstrass. Esquissons une manière de démontrer la lissité.
Considérons le sous-schéma en groupes
$$
H =
\left\{
\left(
\begin{matrix}
u^2 & s & 0 \\
0 & u^3 & 0 \\
r & t & 1
\end{matrix}
\right)
\middle|
\begin{matrix}
u\text{ inversible} \\
s, r, t\text{ arbitraires}
\end{matrix}
\right\}
\subset
\text{GL}_{3, \mathbf{Z}}
$$
Il existe une action $H \times W \to W$ de $H$ sur le schéma de Weierstrass
$W$. Pour en trouver les équations, il suffit d'expliciter l'effet, sur
l'équation générale de Weierstrass, d'un changement de coordonnées donné par
une matrice de $H$. On constate alors que les assertions suivantes sont vraies :
\begin{enumerate}
\item toute courbe elliptique $(E, f, 0)/S$ admet, localement pour la topologie
de Zariski sur $S$, une équation de Weierstrass ;
\item deux équations de Weierstrass quelconques de $(E, f, 0)$ diffèrent,
localement pour la topologie de Zariski, par un élément de $H$.
\end{enumerate}
En considérant le produit fibré
$S \times_{\mathcal{M}_{1, 1}} W =
\mathit{Isom}_{S \times W}(E \times W, S \times E_W)$,
nous en concluons que le morphisme $W \to \mathcal{M}_{1, 1}$ est un
$H$-torseur. Puisque $H \to \Spec(\mathbf{Z})$ est lisse et que les torseurs
sous des schémas en groupes lisses sont lisses, le résultat s'ensuit.
\end{proof}

\begin{remark}
\label{remark-quotient-stack}
L'argument esquissé ci-dessus montre en fait que
$\mathcal{M}_{1, 1} = [W/H]$ est un champ quotient global.
Dans environ 50\% des cas, un argument établissant qu'un champ de modules est
algébrique montre aussi qu'il s'agit d'un champ quotient global.
\end{remark}



\section{Propriétés des champs algébriques}
\label{section-properties}

\noindent
Nous savons maintenant que $\mathcal{M}_{1, 1}$ est un champ algébrique.
Que pouvons-nous en faire ? Ce n'est pas tant le fait qu'il soit algébrique
qui nous aide ici que le point de vue selon lequel les propriétés de
$\mathcal{M}_{1, 1}$ doivent être encodées dans celles des morphismes
$S \to \mathcal{M}_{1, 1}$, c'est-à-dire dans les familles de courbes
elliptiques. En voici quelques exemples.

\medskip\noindent
{\bf Propriétés locales :}
$$
\mathcal{M}_{1, 1} \to \Spec(\mathbf{Z})\text{ est lisse}
\Leftrightarrow
W \to \Spec(\mathbf{Z})\text{ est lisse}
$$
{\bf Idée}. Les propriétés locales d'un champ algébrique sont encodées dans
les propriétés locales de son recouvrement lisse.

\medskip\noindent
{\bf Propriétés globales :}
$$
\begin{matrix}
\mathcal{M}_{1, 1}\text{ est quasi-compact} \Leftarrow W\text{ est quasi-compact}
\\
\mathcal{M}_{1, 1}\text{ est irréductible} \Leftarrow W\text{ est irréductible}
\end{matrix}
$$
{\bf Idée}. Certaines propriétés globales d'un champ algébrique se lisent sur
la propriété correspondante d'un recouvrement lisse
{\it convenable}\footnote{Je suppose qu'il peut exister un champ algébrique
irréductible dépourvu de recouvrement lisse irréductible ; mais, si tel est le
cas, il doit être assez pathologique !}.

\medskip\noindent
{\bf Faisceaux quasi-cohérents :}
$$
\QCoh(\mathcal{O}_{\mathcal{M}_{1, 1}}) =
H\text{-modules quasi-cohérents équivariants sur }W
$$
{\bf Idée}. D'une part, un module quasi-cohérent sur
$\mathcal{M}_{1, 1}$ devrait correspondre à un faisceau quasi-cohérent
$\mathcal{F}_{S, e}$ sur $S$ pour chaque morphisme
$e : S \to \mathcal{M}_{1, 1}$. Cela vaut en particulier pour le morphisme
$(E_W, f_W, 0_W) :  W \to \mathcal{M}_{1, 1}$. Comme ce morphisme est
$H$-équivariant, le module quasi-cohérent $\mathcal{F}_W$ ainsi obtenu est
$H$-équivariant. Réciproquement, à partir d'un module $H$-équivariant, on peut
retrouver les faisceaux $\mathcal{F}_{S, e}$ par la théorie de la descente, en
partant de l'observation que $S \times_{e, \mathcal{M}_{1, 1}} W$ est un
$H$-torseur.

\medskip\noindent
{\bf Groupe de Picard :}
$$
\Pic(\mathcal{M}_{1, 1}) = \Pic_H(W) = \mathbf{Z}/12\mathbf{Z}
$$
{\bf Idée}. Nous avons vu ci-dessus la première égalité. Remarquons que
$\Pic(W) = 0$, car l'anneau
$\mathbf{Z}[a_1, a_2, a_3, a_4, a_6, 1/\Delta]$
a un groupe des classes trivial. Il existe une suite exacte
$$
\mathbf{Z}\Delta \to
\Pic_H(\mathbf{A}^5_{\mathbf{Z}}) \to
\Pic_H(W) \to 0
$$
Le groupe du milieu est égal à $\Hom(H, \mathbf{G}_m) = \mathbf{Z}$.
L'image de $\Delta$ est $12$, car $\Delta$ est de degré $12$.
Cet argument est essentiellement correct ; voir \cite{PicM11}.

\medskip\noindent
{\bf Cohomologie \'etale :} soit $\Lambda$ un anneau.
Il existe une suite spectrale du premier quadrant convergeant vers
$H^{p + q}_\etale(\mathcal{M}_{1, 1}, \Lambda)$, dont la page $E_2$ est
$$
E_2^{p, q} = H_\etale^q(W \times H \times \ldots \times H, \Lambda)
\quad(p\text{ facteurs }H)
$$
{\bf Idée}. Remarquons que
$$
W \times_{\mathcal{M}_{1, 1}} W \times_{\mathcal{M}_{1, 1}}
\ldots \times_{\mathcal{M}_{1, 1}} W
= W \times H \times \ldots \times H
$$
parce que $W \to \mathcal{M}_{1, 1}$ est un $H$-torseur. Cette suite
spectrale est la suite spectrale de {\v C}ech vers la cohomologie associée au
recouvrement lisse $\{W \to \mathcal{M}_{1, 1}\}$. Par exemple,
$H^0_\etale(\mathcal{M}_{1, 1}, \Lambda) = \Lambda$ parce que
$W$ est connexe, et $H^1_\etale(\mathcal{M}_{1, 1}, \Lambda) = 0$ parce que
$H^1_\etale(W, \Lambda) = 0$ (cela exige évidemment une démonstration).
Bien entendu, le recouvrement lisse $W \to \mathcal{M}_{1, 1}$ n'est
peut-être pas « optimal » pour calculer la cohomologie \'etale.














\begin{multicols}{2}[\section{Autres chapitres}]
\noindent
Préliminaires
\begin{enumerate}
\item \hyperref[introduction-section-phantom]{Introduction}
\item \hyperref[conventions-section-phantom]{Conventions}
\item \hyperref[sets-section-phantom]{Théorie des ensembles}
\item \hyperref[categories-section-phantom]{Catégories}
\item \hyperref[topology-section-phantom]{Topologie}
\item \hyperref[sheaves-section-phantom]{Faisceaux sur les espaces}
\item \hyperref[sites-section-phantom]{Sites et faisceaux}
\item \hyperref[stacks-section-phantom]{Champs}
\item \hyperref[fields-section-phantom]{Corps}
\item \hyperref[algebra-section-phantom]{Algèbre commutative}
\item \hyperref[brauer-section-phantom]{Groupes de Brauer}
\item \hyperref[homology-section-phantom]{Algèbre homologique}
\item \hyperref[derived-section-phantom]{Catégories dérivées}
\item \hyperref[simplicial-section-phantom]{Méthodes simpliciales}
\item \hyperref[more-algebra-section-phantom]{Compléments d'algèbre}
\item \hyperref[smoothing-section-phantom]{Lissage des morphismes d'anneaux}
\item \hyperref[modules-section-phantom]{Faisceaux de modules}
\item \hyperref[sites-modules-section-phantom]{Modules sur les sites}
\item \hyperref[injectives-section-phantom]{Objets injectifs}
\item \hyperref[cohomology-section-phantom]{Cohomologie des faisceaux}
\item \hyperref[sites-cohomology-section-phantom]{Cohomologie sur les sites}
\item \hyperref[dga-section-phantom]{Algèbre différentielle graduée}
\item \hyperref[dpa-section-phantom]{Algèbre à puissances divisées}
\item \hyperref[sdga-section-phantom]{Faisceaux différentiels gradués}
\item \hyperref[hypercovering-section-phantom]{Hyperrecouvrements}
\end{enumerate}
Schémas
\begin{enumerate}
\setcounter{enumi}{25}
\item \hyperref[schemes-section-phantom]{Schémas}
\item \hyperref[constructions-section-phantom]{Constructions de schémas}
\item \hyperref[properties-section-phantom]{Propriétés des schémas}
\item \hyperref[morphisms-section-phantom]{Morphismes de schémas}
\item \hyperref[coherent-section-phantom]{Cohomologie des schémas}
\item \hyperref[divisors-section-phantom]{Diviseurs}
\item \hyperref[limits-section-phantom]{Limites de schémas}
\item \hyperref[varieties-section-phantom]{Variétés}
\item \hyperref[topologies-section-phantom]{Topologies sur les schémas}
\item \hyperref[descent-section-phantom]{Descente}
\item \hyperref[perfect-section-phantom]{Catégories dérivées de schémas}
\item \hyperref[more-morphisms-section-phantom]{Compléments sur les morphismes}
\item \hyperref[flat-section-phantom]{Compléments sur la platitude}
\item \hyperref[groupoids-section-phantom]{Schémas en groupoïdes}
\item \hyperref[more-groupoids-section-phantom]{Compléments sur les schémas en groupoïdes}
\item \hyperref[etale-section-phantom]{Morphismes étales de schémas}
\end{enumerate}
Thèmes de théorie des schémas
\begin{enumerate}
\setcounter{enumi}{41}
\item \hyperref[chow-section-phantom]{Homologie de Chow}
\item \hyperref[intersection-section-phantom]{Théorie de l'intersection}
\item \hyperref[pic-section-phantom]{Schémas de Picard des courbes}
\item \hyperref[weil-section-phantom]{Théories cohomologiques de Weil}
\item \hyperref[adequate-section-phantom]{Modules adéquats}
\item \hyperref[dualizing-section-phantom]{Complexes dualisants}
\item \hyperref[duality-section-phantom]{Dualité pour les schémas}
\item \hyperref[discriminant-section-phantom]{Discriminants et différentes}
\item \hyperref[derham-section-phantom]{Cohomologie de de Rham}
\item \hyperref[local-cohomology-section-phantom]{Cohomologie locale}
\item \hyperref[algebraization-section-phantom]{Géométrie algébrique et formelle}
\item \hyperref[curves-section-phantom]{Courbes algébriques}
\item \hyperref[resolve-section-phantom]{Résolution des surfaces}
\item \hyperref[models-section-phantom]{Réduction semi-stable}
\item \hyperref[functors-section-phantom]{Foncteurs et morphismes}
\item \hyperref[equiv-section-phantom]{Catégories dérivées de variétés}
\item \hyperref[pione-section-phantom]{Groupes fondamentaux de schémas}
\item \hyperref[etale-cohomology-section-phantom]{Cohomologie étale}
\item \hyperref[crystalline-section-phantom]{Cohomologie cristalline}
\item \hyperref[proetale-section-phantom]{Cohomologie pro-étale}
\item \hyperref[relative-cycles-section-phantom]{Cycles relatifs}
\item \hyperref[more-etale-section-phantom]{Compléments de cohomologie étale}
\item \hyperref[trace-section-phantom]{La formule des traces}
\end{enumerate}
Espaces algébriques
\begin{enumerate}
\setcounter{enumi}{64}
\item \hyperref[spaces-section-phantom]{Espaces algébriques}
\item \hyperref[spaces-properties-section-phantom]{Propriétés des espaces algébriques}
\item \hyperref[spaces-morphisms-section-phantom]{Morphismes d'espaces algébriques}
\item \hyperref[decent-spaces-section-phantom]{Espaces algébriques décents}
\item \hyperref[spaces-cohomology-section-phantom]{Cohomologie des espaces algébriques}
\item \hyperref[spaces-limits-section-phantom]{Limites d'espaces algébriques}
\item \hyperref[spaces-divisors-section-phantom]{Diviseurs sur les espaces algébriques}
\item \hyperref[spaces-over-fields-section-phantom]{Espaces algébriques sur des corps}
\item \hyperref[spaces-topologies-section-phantom]{Topologies sur les espaces algébriques}
\item \hyperref[spaces-descent-section-phantom]{Descente et espaces algébriques}
\item \hyperref[spaces-perfect-section-phantom]{Catégories dérivées d'espaces}
\item \hyperref[spaces-more-morphisms-section-phantom]{Compléments sur les morphismes d'espaces}
\item \hyperref[spaces-flat-section-phantom]{Platitude sur les espaces algébriques}
\item \hyperref[spaces-groupoids-section-phantom]{Groupoïdes dans les espaces algébriques}
\item \hyperref[spaces-more-groupoids-section-phantom]{Compléments sur les groupoïdes dans les espaces}
\item \hyperref[bootstrap-section-phantom]{Amorçage}
\item \hyperref[spaces-pushouts-section-phantom]{Sommes amalgamées d'espaces algébriques}
\end{enumerate}
Thèmes de géométrie
\begin{enumerate}
\setcounter{enumi}{81}
\item \hyperref[spaces-chow-section-phantom]{Groupes de Chow des espaces}
\item \hyperref[groupoids-quotients-section-phantom]{Quotients de groupoïdes}
\item \hyperref[spaces-more-cohomology-section-phantom]{Compléments sur la cohomologie des espaces}
\item \hyperref[spaces-simplicial-section-phantom]{Espaces simpliciaux}
\item \hyperref[spaces-duality-section-phantom]{Dualité pour les espaces}
\item \hyperref[formal-spaces-section-phantom]{Espaces algébriques formels}
\item \hyperref[restricted-section-phantom]{Algébrisation des espaces formels}
\item \hyperref[spaces-resolve-section-phantom]{Résolution des surfaces revisitée}
\end{enumerate}
Théorie des déformations
\begin{enumerate}
\setcounter{enumi}{89}
\item \hyperref[formal-defos-section-phantom]{Théorie formelle des déformations}
\item \hyperref[defos-section-phantom]{Théorie des déformations}
\item \hyperref[cotangent-section-phantom]{Le complexe cotangent}
\item \hyperref[examples-defos-section-phantom]{Problèmes de déformation}
\end{enumerate}
Champs algébriques
\begin{enumerate}
\setcounter{enumi}{93}
\item \hyperref[algebraic-section-phantom]{Champs algébriques}
\item \hyperref[examples-stacks-section-phantom]{Exemples de champs}
\item \hyperref[stacks-sheaves-section-phantom]{Faisceaux sur les champs algébriques}
\item \hyperref[criteria-section-phantom]{Critères de représentabilité}
\item \hyperref[artin-section-phantom]{Axiomes d'Artin}
\item \hyperref[quot-section-phantom]{Espaces de Quot et de Hilbert}
\item \hyperref[stacks-properties-section-phantom]{Propriétés des champs algébriques}
\item \hyperref[stacks-morphisms-section-phantom]{Morphismes de champs algébriques}
\item \hyperref[stacks-limits-section-phantom]{Limites de champs algébriques}
\item \hyperref[stacks-cohomology-section-phantom]{Cohomologie des champs algébriques}
\item \hyperref[stacks-perfect-section-phantom]{Catégories dérivées de champs}
\item \hyperref[stacks-introduction-section-phantom]{Introduction aux champs algébriques}
\item \hyperref[stacks-more-morphisms-section-phantom]{Compléments sur les morphismes de champs}
\item \hyperref[stacks-geometry-section-phantom]{La géométrie des champs}
\end{enumerate}
Thèmes de théorie des espaces de modules
\begin{enumerate}
\setcounter{enumi}{107}
\item \hyperref[moduli-section-phantom]{Champs de modules}
\item \hyperref[moduli-curves-section-phantom]{Espaces de modules de courbes}
\end{enumerate}
Divers
\begin{enumerate}
\setcounter{enumi}{109}
\item \hyperref[examples-section-phantom]{Exemples}
\item \hyperref[exercises-section-phantom]{Exercices}
\item \hyperref[guide-section-phantom]{Guide bibliographique}
\item \hyperref[desirables-section-phantom]{Desiderata}
\item \hyperref[coding-section-phantom]{Style de programmation}
\item \hyperref[obsolete-section-phantom]{Obsolète}
\item \hyperref[fdl-section-phantom]{Licence de documentation libre GNU}
\item \hyperref[index-section-phantom]{Index généré automatiquement}
\end{enumerate}
\end{multicols}


\bibliography{my}
\bibliographystyle{amsalpha}

\end{document}
